\chapter{同调}

\section{自由Abel群}

Abel群$A$中的元素$x$有逆元$-x$，
整数$n \in \mathbb Z$给出了

$$
nx = \begin{cases}
  \underbrace{x + \cdots + x}_n & (n > 0) \\
  0 & (n = 0) \\
  \underbrace{(-x) + \cdots + (-x)}_{|n|} & (n > 0) \\
\end{cases}
$$

进而取子集$\{x_i\} \subseteq A$构造和：

$$
\sum_i n_i x_i
$$

这样形式的全体构成了$\{x_i\}$生成的子群。
若生成集合有限则称有限生成群，
若$\sum_i n_i x_i = 0$的条件是所有$n_i = 0$退化，
称生成集$\{x_i\}$线性无关。
由$r$个线性无关元有限生成的群称为秩$r$自由Abel群。
$\mathbb Z$是$1$或$-1$生成的自由Abel群。


\section{正合}

环$\Lambda \in \textbf{Ring}$上的左模范畴记$\Lambda-\textbf{Mod}$，
考虑模同态$f \in \Lambda-\textbf{Mod}(X,Y)$，
产生了同态的像与核：

$$
\text{Im} \ f = \{ y \in Y \mid \exists x \in X,\ f(x) = y\},\ 
\quad
\text{Ker} \ f = \{x\in X \mid f(x) < 0\}
$$

它们都是子模，
从而可以用商模定义：

$$
\text{Coim} \ f = X/\text{Ker} \ f,\ 
\quad
\text{Coker} \ f = Y/\text{Im} \ f
$$

由$f$有同态：

$$
\begin{aligned}
  f^\prime : \text{Coim} \ f = X/\text{Ker} \ f &\to Y \\
  x + \text{Ker} \ f &\mapsto f^\prime(x + \text{Ker} \ f) = f(x)
\end{aligned}
$$

显然$\text{Im} \ f^\prime = \text{Im} \ f$，
$f^\prime$是满同态。
$f^\prime(x + \text{Ker} \ f) = f(x) = 0$时，
$x \in \text{Ker} \ f$，
即$x + \text{Ker} \ f = 0$，
有$\text{Ker} \ f^\prime = 0$，
$f^\prime$是单同态。
$\text{Coim} \ f$描述的就是同构$ = X/\text{Ker} \ f \simeq \text{Im} \ f$。
当$\text{Ker} \ f = 0$时称为单同态，
$\text{Im} \ f = Y$时称为满同态，
单满同态为同构，
此时

$$
\text{Coim} \ f = X/\text{Ker} \ f \simeq X \simeq \text{Im} \ f = Y
$$

像和核都是子模，
若将二者放在一个模中比较，
即两个首尾相接的同态中，
$\text{Im} \ f$和$\text{Ker} \ g$都是$Y$的子模：

$$
X \stackrel{f}{\longrightarrow} Y \stackrel{g}{\longrightarrow} Z
$$

若复合映射$X \stackrel{g \circ f}{\longrightarrow} Z$
退化为$g \circ f = 0$，
此时$\text{Im} \ f \subseteq \text{Ker} \ g$。
当$\text{Im} \ f = \text{Ker} \ g$时称为正合。
正合可以用于描述单/满同态：

$$
0 \stackrel{i}{\longrightarrow} X \stackrel{f}{\longrightarrow} Y,\
\quad
Y \stackrel{g}{\longrightarrow} Z \stackrel{p}{\longrightarrow} 0
$$

以上两个正合意味着$\text{Im} \ i = \text{Ker} \ f = 0$，
它说明$f$是单同态；
以及$\text{Im} \ g = \text{Ker} \ p = Z$，
它说明$g$是满同态。
将二者合并得到短正合序列，
描述了单同态$f$和满同态$g$在$Y$处正合：

$$
0 \longrightarrow X \stackrel{f}{\longrightarrow} Y \stackrel{g}{\longrightarrow} Z \longrightarrow 0
$$

单同态$f$把$X$以同构的意义嵌入$Y$作为子模，
在自然投射$\pi: Y \to Y/X$中，
子模$X \simeq \text{Ker} \ \pi$，
于是：

$$
0 \longrightarrow X \stackrel{f}{\longrightarrow} Y \stackrel{\pi}{\longrightarrow} Y/X \longrightarrow 0
$$


\section{复形}


\section{同调群}



